\documentclass[11pt,a4paper]{article}
\usepackage[T1]{fontenc}
\usepackage{lmodern}
\usepackage[margin=25mm]{geometry}
\usepackage{amsmath,amssymb,booktabs,microtype}
\usepackage[hidelinks]{hyperref}
\usepackage{enumitem}
\setlist{nosep,leftmargin=*}
\setlength{\parskip}{4pt}
\setlength{\parindent}{1em}
\setlength{\emergencystretch}{2em}
\clubpenalty=10000
\widowpenalty=10000
\hypersetup{pdftitle={ConnectCoin: A Peer-to-Peer Cryptocurrency for TLS-Based Rewards},pdfauthor={Papaulo},
  pdfsubject={A Peer-to-Peer Cryptocurrency for TLS-Based Rewards}}
\newcommand{\CONN}{\mathrm{CONN}}
\newcommand{\SHA}{\operatorname{SHA256}}
\newcommand{\TagHash}{\operatorname{TaggedHash}}
\newcommand{\txid}{\operatorname{txid}}
\newcommand{\LE}{\operatorname{LE32}}
\newcommand{\concat}{\mathbin{\|}}
\title{ConnectCoin\\[3pt]\large A Peer-to-Peer Cryptocurrency for TLS-Based Rewards}
\author{Papaulo}
\date{September 2026}

\begin{document}
\begin{center}
{\LARGE ConnectCoin\par}
\vspace{5pt}
{\large A Peer-to-Peer Cryptocurrency for TLS-Based Rewards\par}
\vspace{8pt}
Papaulo\par
\vspace{8pt}
September 2026
\end{center}
\vspace{6pt}

\begin{abstract}
ConnectCoin is a Bitcoin-derived blockchain that associates transaction outputs
with cryptographically verifiable TLS interactions. A participant creates a
Pay-to-Connect output specifying a domain, a reward, a connection-work
target, and permitted signature schemes. Another participant prepares a spending transaction and obtains a
server-authenticated TLS handshake transcript bound to that transaction. Nodes
verify the submitted evidence without contacting the domain. A hash target
selects which transcripts qualify for payment, allowing fewer proofs to be
recorded than the number of candidate transcripts examined. The work hash
covers the authenticated handshake prefix, with the server signature retained
as evidence. RandomX proof-of-work separately secures block
ordering. A declining issuance schedule and a subsidy reduction proportional
to transaction weight establish the monetary and block-space incentives.
This paper describes the construction and its economic policy.
\end{abstract}

\section{Introduction}

A digital payment can reward a network interaction only if the parties can agree
on evidence that the interaction occurred. A recipient's own report is not
sufficient evidence, while asking every blockchain node to repeat the interaction
would make validation depend on network access, changing server state, and local
timing. The objective is therefore to separate obtaining an external response
from verifying the evidence used to claim a payment.

TLS provides a useful starting point: a server authenticates part of its
handshake using a certificate-backed signing key. ConnectCoin makes a restricted
form of that authenticated transcript a transaction-spending condition. A claim
contains the messages and certificates needed for deterministic verification,
and the server's response is bound to the transaction that receives the reward.
Payment does not require the funder to approve the claimant after seeing the
result.

The proposal has two layers. \emph{Proof-of-Connect} defines the evidence and
its connection-work condition. \emph{Pay-to-Connect} attaches that condition to
an unspent transaction output. These are distinct from block production:
RandomX proof-of-work orders transactions, whereas P2C proofs spend rewards
already supplied by other participants. A successful P2C claim does not itself
mint currency or create a block.

The intended result is a market for narrowly defined, verifiable interactions,
not a universal oracle for Internet activity. A transcript does not establish
that a human visited a website, that content was delivered, or that the server
is independent of the claimant.

\section{Network Specifications}

ConnectCoin reuses the Bitcoin Core codebase and the UTXO model, in which inputs
consume previous outputs and new outputs define subsequent spending rights.
Transactions are ordered in blocks, and nodes select the valid chain with the
greatest accumulated proof-of-work~\cite{bitcoin}. Changes to transaction encoding,
proof-of-work, and network parameters make ConnectCoin incompatible with the
Bitcoin network. Despite the fact that partial implementation was reused, there will be no compatibility with
Bitcoin balances, wallets, or historical blocks.

\subsection*{Block production and capacity}

Block proof-of-work uses RandomX V2 over the serialized 80-byte header. RandomX
is designed around general-purpose CPU execution and substantial memory
requirements~\cite{randomx}. By favoring general-purpose CPUs over GPUs and seeking
to minimize the efficiency advantage of application-specific integrated circuits
(ASICs), this design aims to support home mining on widely available hardware and
encourage a more decentralized distribution of mining power.

The configured target interval is 10 seconds. The ordinary difficulty adjustment
period is 8,640 blocks, corresponding to one day at that interval. Difficulty
adjustment must not be confused with the RandomX key schedule: the latter uses
12,288-block epochs with a 64-block lag and a network-specific bootstrap key.

\begin{center}
\begin{tabular}{@{}ll@{}}
\toprule
Parameter & Public test-network setting \\
\midrule
Target block interval & 10 seconds \\
Ordinary difficulty adjustment & 8,640 blocks \\
Maximum block weight & 50,000,000 weight units \\
Coinbase maturity & 100 blocks \\
Initial recurring subsidy & 15 CONN per block \\
Halving interval & 3,000,000 blocks \\
Smallest monetary unit & $10^{-10}$ CONN, one connect \\
\bottomrule
\end{tabular}
\end{center}

Weight follows the distinction between transaction base data and witness data
used by segregated witness~\cite{segwit}. If $b$ is the stripped size and $t$ is
the total serialized size, then
\begin{equation}
  w=3b+t=4b+(t-b).
\end{equation}
Consequently, 50 million weight units are not generally 50 million bytes of
ordinary transaction data. Depending on the proportion of witness data,
serialized blocks can approach 50 MB within this weight limit.

\subsection*{Typed outputs}

Instead of general transaction-output scripts, the current protocol recognizes
two output types. Type 1 contains a 32-byte x-only secp256k1 public key and is
spent with a 64-byte BIP340 Schnorr signature~\cite{schnorr}. Its transaction
signature digest follows the BIP341-style default signing construction
\cite{taproot}. Type 2 is the Pay-to-Connect output defined in Section~4.
Unknown output types are rejected. This restricted format intentionally does
not implement arbitrary Bitcoin Script spending conditions.

\section{Proof-of-Connect}

\subsection*{Binding a response to a claim}

Let $C$ be a prepared claim transaction, with the reward destination and all
non-witness fields already fixed, and let $i$ be its P2C input index. The
connection challenge is
\begin{equation}
  c=\TagHash_{\mathrm{claim}}\bigl(\txid(C)\concat\LE(i)\bigr).
  \label{eq:challenge}
\end{equation}
The tag is \texttt{ConnectCoin/P2C/claim/v1}. Transaction-ID bytes use the
protocol serialization, and the input index is a four-byte little-endian
integer. Tagged hashing uses the domain-separation construction
$\SHA(\SHA(\mathrm{tag})\concat\SHA(\mathrm{tag})\concat m)$~\cite{schnorr}.

The claimant puts the exact 32 bytes of $c$ in \texttt{ClientHello.random}
and sets the requested domain in the server-name extension. A new connection
attempt can retain this challenge while generating a fresh key share. The
server need not run a ConnectCoin-specific application: a compatible TLS
implementation authenticates the handshake containing these client fields.

The transaction ID excludes witness data, avoiding a circular dependency
between the claim and its proof. It nevertheless commits to the consumed
outputs and the payout transaction. Changing a recipient, fee, or other
non-witness field changes the challenge and requires new server-authenticated
evidence.

\subsection*{Authenticated evidence}

Proof version 2 contains the raw handshake messages
\begin{equation}
  \tau=CH\concat SH\concat EE\concat CRT\concat CV,
  \label{eq:transcript}
\end{equation}
Here $CH$ is \texttt{ClientHello}, $SH$ is \texttt{ServerHello},
$EE$ is \texttt{EncryptedExtensions}, $CRT$ is \texttt{Certificate},
and $CV$ is \texttt{CertificateVerify}. Each includes its TLS handshake header.
The proof is the byte \texttt{0x02} followed by these messages, in this order.
It excludes TLS record framing, session secrets, \texttt{Finished}, and
application data~\cite{connectcoin}.

The TLS 1.3 server signature authenticates the handshake prefix through
\texttt{Certificate}, using the TLS server context and the public key from
the leaf certificate~\cite{tls}. A validator checks that the certificate chain
is accepted for the specified domain, that the expected challenge is present,
and that the server signature verifies. The strict inference is therefore an
authenticated transcript under a key authorized for that domain, not proof
that every phase of a TLS session completed.

ConnectCoin accepts a deliberately restricted TLS 1.3 profile: AES-128-GCM
or ChaCha20-Poly1305 with SHA-256, X25519 or secp256r1 key shares, and
ECDSA P-256/SHA-256 or either supported RSA-PSS/SHA-256 authentication scheme.
The selected signature scheme must be offered in the ClientHello and permitted
by the output's mask. RSA keys must contain at least 2,048 bits; RSA-PSS uses
MGF1-SHA-256 and a 32-byte salt. RSA encryption keys and restricted RSA-PSS
keys are distinguished and checked against the selected scheme.
Resumption, early data, and HelloRetryRequest are excluded. This
reduces the set of states that consensus must interpret, at the cost of
compatibility with some servers. The exact versioned profile, rather than
the entirety of any TLS standard, determines valid P2C evidence. Later TLS
specification revisions do not automatically change consensus~\cite{tlsupdate}.

Certificate roots form an immutable, versioned consensus bundle. Validators
use that bundle rather than their operating system's changing trust store.
The validity dates of supplied certificates are evaluated at the previous block's median time past
when validating a block, and at the current tip's median time past for mempool
admission. No DNS lookup or live certificate-status request is part of this
verification.

\subsection*{Connection-work selection}

Define the authenticated prefix $M=CH\concat SH\concat EE\concat CRT$.
The candidate work hash is
\begin{equation}
  H=\TagHash_{\mathrm{work}}(M), \qquad H\leq T,
  \label{eq:work}
\end{equation}
where the tag is \texttt{ConnectCoin/P2C/work/v2} and $T$ is the output's
256-bit target. These four messages include their handshake headers.
The entire \texttt{CertificateVerify} message and the proof-version byte are
excluded from this hash. The digest is interpreted using ConnectCoin's little-endian
256-bit integer representation. The target comparison is inclusive.

The server's CertificateVerify signature authenticates $M$ and is required
for a valid proof.

For independent, uniformly distributed candidate hashes, the success
probability and expected number of candidate evaluations are
\begin{equation}
  p(T)=\frac{T+1}{2^{256}}, \qquad \mathbb{E}[K]=\frac{1}{p(T)}.
  \label{eq:probability}
\end{equation}
The probability of seeing no qualifying candidate after $n$ evaluations is
$(1-p)^n$. For example, $T=2^{246}-1$ accepts hashes with ten leading zero
bits in their numeric representation and gives $p=1/1024$.

\section{Pay-to-Connect}

\subsection*{Funding and redemption}

A Pay-to-Connect output has a value $V$ and payload
\begin{equation}
  (d,T,r,m),
\end{equation}
where $d$ is a lower-case ASCII domain satisfying the protocol's label grammar,
$T$ is its connection-work target, $r$ identifies the trusted-root bundle, and
$m$ is a mandatory one-byte \texttt{signature\_algorithms\_mask}.
Domain length is encoded in one byte, the target occupies 32 bytes, and the
root version occupies four. With the eight-byte amount and one-byte output
type, the complete serialized output occupies $47+|d|$ bytes.

\begin{center}
\begin{tabular}{@{}cll@{}}
\toprule
Mask value & CertificateVerify scheme & TLS identifier \\
\midrule
1 & \texttt{ecdsa\_secp256r1\_sha256} & \texttt{0x0403} \\
2 & \texttt{rsa\_pss\_rsae\_sha256} & \texttt{0x0804} \\
4 & \texttt{rsa\_pss\_pss\_sha256} & \texttt{0x0809} \\
\bottomrule
\end{tabular}
\end{center}

For example, mask 7 permits all three schemes; mask 6 permits both RSA schemes. Zero and
reserved bits are invalid. This mask restricts server authentication, not
the issuer signatures along the certificate chain. The RPC default is 7.
When funding through the graphical wallet, a capability probe offers only
the two RSA schemes during the three-second confirmation delay. A full,
authenticated handshake completed before the deadline selects 6; failure or
timeout retains 7. The choice freezes before confirmation is enabled.
This advisory probe uses local time and the selected endpoint's current
capability; it cannot guarantee future acceptance under chain-time validation.
One advantage of this preference is the lower computational cost of RSA
signature verification.

Anyone can fund this output, and anyone satisfying its spending condition can
claim it. There is no creator signature required at redemption and no
creator-only refund branch. In particular, the current type does not contain
a deadline or a mechanism to recover funds if the domain becomes permanently
unusable. Funding a bounty is a commitment to its encoded conditions.

\begin{figure}[ht]
\centering
\fbox{\shortstack{Funder creates\\P2C output}}
$\longrightarrow$
\fbox{\shortstack{Claimant fixes\\claim transaction}}
$\longrightarrow$
\fbox{\shortstack{Obtain qualifying\\TLS evidence}}
$\longrightarrow$
\fbox{\shortstack{Nodes verify\\and settle}}
\caption{Only the claimant obtains the external response. Validators use the
evidence carried by the claim transaction.}
\end{figure}

A simple claim spends one bounty and creates an ordinary output paying
$V-f$ to the claimant, where $f$ is the transaction fee.

% Keep this validation paragraph intact when the preceding figure floats.
\begin{samepage}
For each P2C input, validators require an existing, spendable output; the
correct proof version and message structure; the transaction-bound challenge;
the matching domain and accepted certificate path; a valid, mask-permitted server signature;
and a work hash satisfying the output target. These checks supplement ordinary
transaction, block, and double-spend validation. Each P2C input has an empty
scriptSig and exactly one proof witness element. The full proof is limited to
64 KiB (65,536 bytes). The Certificate message is
limited to 48 KiB, containing at most eight certificates of at most 16 KiB
each; CertificateVerify has its own 8 KiB limit.
\par
\end{samepage}

\subsection*{Competition and reuse}

Each bounty is a single-use UTXO. Multiple claimants may search concurrently,
but only one conflicting spend can survive in the selected chain. Producing
a valid proof does not guarantee payment: another claim can be confirmed
first, and a reorganization can change the accepted history. A third party
can relay a completed claim but cannot redirect its outputs without changing
the challenge in Equation~\eqref{eq:challenge}.

The input-index binding also prevents simply copying one proof to several
inputs of the same transaction. A funding transaction can create many
independent rewards for the same domain, but each output remains a separate
claim opportunity. Splitting a fixed budget can provide more independently
settled outcomes while increasing funding and claim overhead; it does not
create additional value.

There is no special P2C confirmation-age rule in consensus for ordinary
outputs beyond the usual spending and lock conditions. The automatic wallet
waits for confirmed bounties as a local policy. Coinbase outputs separately
retain their 100-block maturity. Search limits, domain scheduling, and
decisions to abandon an unproductive bounty do not invalidate otherwise
valid claims.

\subsection*{Automatic claim selection}

The wallet first discovers eligible outputs using an in-memory P2C catalog
initialized from the UTXO set. Later updates apply newly connected blocks and
undo reorganized blocks; unavailable history can require rebuilding the
catalog. The discovery window is configurable: by default it covers the latest
600 blocks, including the tip, approximately 100 minutes at the target interval.
Unlimited history includes all confirmed ages (\texttt{recent\_blocks}=0).
The wallet warns that large windows can include old, unproductive bounties.
This is a soft discovery window, not on-chain expiry: crossing its boundary
does not cancel an in-flight attempt.

The wallet refreshes the candidate set every five seconds and selects a domain
for each new connection. Assignments alternate a fair turn among eligible domains and a turn
for the domain with the highest estimated reward rate. Thus rewards influence
cross-domain priority without making it absolute. The economic score uses
the best eligible bounty within a domain, not the sum of its advertised rewards.
There is no fixed per-domain connection quota.

% Keep the estimator's definition with its equation.
\begin{samepage}
For each pair of domain and signature mask, the wallet retains the last 100
completed TCP/TLS attempts. Write $s$ for captures reaching CertificateVerify
and $t$ for the sum of attempt durations in seconds. Its smoothed capture rate is
\begin{equation}
 q=\frac{0.1+s}{0.02+t}.
 \label{eq:rate}
\end{equation}
\par
\end{samepage}
Before any samples, $q=5$ captures per second of connection effort. Failures
and timeouts contribute time without success; locally canceled incomplete
attempts are excluded. DNS resolution, scheduling waits, and later certificate
verification are not timed here. Reaching CertificateVerify is a capture
success, not necessarily a valid or winning claim. Concurrent attempt
durations are added independently, so this is not wall-clock throughput.

With net bounty value $V-f$ measured in connects, the estimated reward rate is
\begin{equation}
 G=(V-f)\,\frac{T+1}{2^{256}}\,q.
 \label{eq:expectedreward}
\end{equation}
Automatic selection requires $G\geq1{,}000$ connects per second of connection
effort. This is a monetary threshold, not 1,000 connections per second.
This estimate does not deduct operating costs or account for competing claims.

The wallet stops starting new attempts for a bounty once its number of
successful captures through CertificateVerify exceeds $2/p(T)$, whether or not
those captures yield valid, winning claims. Failed connections and incomplete
local cancellations do not consume this count. With sequential attempts, the
cutoff is $N=\lfloor2/p(T)\rfloor+1$ captures: five when $p=1/2$.
Concurrent attempts already underway may finish and submit proofs after the
cutoff. Counters remain local to the loaded wallet, surviving stop/start until
their bounty leaves discovery and pending work drains; they reset on wallet
reload. For independent uniform
candidates with small $p$, no win in approximately $2/p$ trials has probability
near $e^{-2}\simeq13.5\%$. This is a resource-management heuristic, not evidence
that the server cheated.

Automatic searching is opt-in. The graphical wallet proposes a limit of 100
connection starts per second and independently defaults to 100 simultaneous
connections. Higher settings and an unlimited rate remain available, with a
prominent warning. HTTPS remains disabled until the user starts it and accepts
the confirmation. The rate is shared by the workers of that wallet, not applied
separately to each domain. Both the block miner and automatic P2C
claims use a receiving destination from the selected wallet by default, while
allowing an explicit supported address. A completed proof is bound to the
destination fixed in its claim; it cannot be retargeted afterward.
All selection rules in this subsection are wallet policy, not additional
consensus conditions for manually submitted claims.

\section{Economic Policy}

\subsection*{Units, allocation, and halvings}

One CONN contains $10^{10}$ connects, the atomic accounting unit. The intended
nominal supply budget is 100 million CONN. The proposed launch allocation is:
\begin{itemize}
 \item \textbf{Development Fund: 5\% (5,000,000 CONN).} Restricted to hosting
 websites and rewarding independent developers of official ConnectCoin
 projects, including ConnectCoin Core, and high-quality unofficial projects
 that contribute substantial value, such as alternative wallets and block
 explorers. We commit to publicly explaining all expenditures
 from the Development Fund.
 \item \textbf{Developer's Fund: 5\% (5,000,000 CONN).} The founder's personal
 allocation, intended to cover personal living expenses and provide an
 incentive for sustained work on ConnectCoin.
 \item \textbf{Ordinary mining: 90\% (90,000,000 CONN).} The nominal recurring
 subsidy budget, subject to integer rounding and the weight-dependent
 reduction described below.
\end{itemize}
The Development Fund and the Developer's Fund are issued together in the
genesis block's coinbase transaction, totaling 10,000,000 CONN.

We commit to keeping cumulative withdrawals from each fund within an
allowance that increases by \textbf{1,000,000 CONN each year}. For each fund
separately, this annual increment is 1\% of the nominal total distribution,
or 20\% of that fund's initial 5,000,000 CONN allocation. It defines a maximum,
not a minimum or a withdrawal target. The annual increment is not recalculated
as 20\% of the remaining balance.

The initial 1,000,000 CONN allowance for each fund becomes available on the
mainnet launch date, with another 1,000,000 CONN added on each calendar anniversary.
Any unused withdrawal allowance carries forward, increasing the amount
available for withdrawal in later annual periods. For example, if no withdrawals
are made from a fund during the first period, up to 2,000,000 CONN may be withdrawn
from that fund during the second period. Each fund retains its own cumulative
allowance and purpose. This schedule authorizes withdrawals only from available
balances; it does not create new allocations after the original funds have
been used.

Let $h\geq1$ be the height of an ordinary block, $L=3{,}000{,}000$ the halving
interval, and $S(h)$ the scheduled subsidy measured in connects. Then
\begin{equation}
 S(h)=\left\lfloor
    \frac{15\cdot10^{10}}{2^{\lfloor h/L\rfloor}}
 \right\rfloor.
 \label{eq:subsidy}
\end{equation}
The initial ordinary subsidy is 15 CONN. Halving occurs by block height, not by
calendar date: $L$ blocks at the target interval take about 347.22 days.
``Approximately annual'' is therefore a description of the target schedule,
not a promise of an exact year between halvings.

\begin{samepage}
With a 10 million CONN initial allocation and every subsequent subsidy fully
claimed, the exact integer schedule gives
\begin{equation}
 10{,}000{,}000+
 \frac{1}{10^{10}}\sum_{h=1}^{\infty}S(h)
 =99{,}999{,}984.9955\ \CONN.
 \label{eq:issuance}
\end{equation}
\end{samepage}
The difference from 100 million follows from integer rounding and the genesis
occupying height zero rather than receiving an additional ordinary subsidy.
Actual issuance can be lower because of weight reductions or unclaimed
subsidies. The separate 100 million CONN amount-range bound is a validation
sanity limit, not itself the issuance formula.

\subsection*{Block-space incentive}

A miner earns the scheduled subsidy for a block containing only its coinbase.
Including other transactions reduces the available subsidy in proportion to
their weight. Let $W=50{,}000{,}000$ weight units denote the
\emph{consensus-enforced maximum total block weight}, and let $w$ be the sum
of all \emph{non-coinbase transaction weights} in the block. For a valid block,
$0\leq w<W$. The reduction in connects is
\begin{equation}
 P(h,w)=\left\lfloor
     \frac{S(h)w}{10W}
 \right\rfloor.
 \label{eq:penalty}
\end{equation}
If $F$ is the sum of transaction fees, the maximum total coinbase output is
\begin{equation}
 R_{\max}(h,w)=F+S(h)-P(h,w).
 \label{eq:reward}
\end{equation}
The header and coinbase count toward the block-weight limit but not toward
$w$. A block containing only its coinbase therefore incurs no reduction,
while increasing non-coinbase weight brings the reduction toward 10\% of the
scheduled subsidy.

At $w=W/2$, the subsidy reduction is 5\%, apart from rounding. Transaction fees can
make total coinbase revenue exceed the base subsidy even when the subsidy
itself has been reduced.

The withheld amount is never created. It is not transferred to a treasury,
burn address, or other participants, and it is not recovered by the issuance
schedule later. P2C claims transfer existing rewards; they are not another
source of monetary emission.

For a package that fits within the block-weight limit, with additional weight
$\Delta w$, the immediate marginal subsidy cost is $P(h,w+\Delta w)-P(h,w)$.
A fee-revenue-oriented miner has an
incentive to include it when its fees exceed that cost. The native block
assembler uses this comparison with locally modified package fees. This is
selection policy, not a consensus requirement that miners maximize revenue.
Arithmetic is performed with exact integer rounding and overflow-safe
intermediates.

\subsection*{Relay fees}

One kilovirtualbyte corresponds to 4,000 weight units. The default relay fee
rate is the next block's subsidy cost for that weight, plus a margin of
1,000 connects per kilovirtualbyte:
\begin{equation}
 r(h)=\left\lfloor\frac{4{,}000\,S(h)}{10W}\right\rfloor+1{,}000
 \quad\text{connects/kvB}.
 \label{eq:relay}
\end{equation}
Initially this is 1,201,000 connects/kvB, or 0.0001201 CONN/kvB; after the first
halving it is 601,000 connects/kvB. The fixed margin remains when the
subsidy-dependent term falls. Mempool conditions and local configuration can
require a different rate. A node's relay policy is not an immutable
minimum-fee rule for consensus-valid blocks.

When fee estimation is unavailable and no nonzero fallback rate is configured,
automatically selected wallet fees use the greatest of the next block's economic
mining floor, the current mempool minimum, and applicable wallet/relay minimums.
Automatic P2C claims reserve a fee budget before fixing the
transaction challenge. Their budget accounts for the maximum permitted proof
size; unused budget remains a fee rather than being returned by changing an
already authenticated payout. Fees therefore also affect eligibility in
Equation~\eqref{eq:expectedreward}.

% Keep the closing argument together; references may follow on the same page.
\clearpage
\section{Conclusion}

We have presented a system for funding TLS-based interactions through open,
predefined spending conditions. We began with the UTXO model, so that every
bounty represents funds already committed and can be redeemed only once.
We then bound each proof to a particular transaction and input, fixing the
recipient before the search begins and preventing the evidence from being
reused for a different input or payout. Compatible servers provide the
required authentication through TLS without running ConnectCoin-specific
software.

With these conditions in place, claimants can compete independently for the
same reward, while the selected chain resolves conflicting spends. The funder
does not need to select a claimant or approve the payment afterward, and
transaction fees can be taken from the bounty itself. Participants choose
which rewards to pursue and how much effort to commit, without making those
choices additional consensus requirements.

Finally, the block-space incentive gives transaction inclusion an explicit
economic cost: miners weigh fee revenue against the subsidy forgone by
including additional transaction weight. Funding, claiming, and block
production remain distinct activities, coordinated through shared transaction
rules rather than individual arrangements between participants.

\begin{thebibliography}{9}
\small\raggedright
\bibitem{bitcoin}
S. Nakamoto. \emph{Bitcoin: A Peer-to-Peer Electronic Cash System}. 2008.
\url{https://bitcoin.org/bitcoin.pdf}.

\bibitem{connectcoin}
ConnectCoin contributors. \emph{ConnectCoin Core}, revision
\texttt{783840b236}, September 2026.
Protocol and monetary rules: \texttt{src/consensus/p2c.cpp},
\texttt{src/consensus/p2c\_x509.cpp}, \texttt{src/kernel/chainparams.cpp},
\texttt{src/validation.cpp}, \texttt{src/wallet/p2c\_worker\_impl.cpp},
\texttt{doc/typed-outputs.md}, and \texttt{doc/p2c-wallet.md}.
\href{https://github.com/connectcoincrypto/connectcoin/tree/783840b2360d6b6f9e8375d5b43dbe10cecf4194}{Source tree and protocol documentation}.

\bibitem{randomx}
tevador and RandomX contributors. \emph{RandomX: Proof of Work Algorithm
Based on Random Code Execution}. Design and reference implementation.
\url{https://github.com/tevador/RandomX}.

\bibitem{segwit}
E. Lombrozo, J. Lau, and P. Wuille. \emph{Segregated Witness (Consensus Layer)}.
BIP 141, 2015. \url{https://github.com/bitcoin/bips/blob/master/bip-0141.mediawiki}.

\bibitem{schnorr}
P. Wuille, J. Nick, and T. Ruffing. \emph{Schnorr Signatures for secp256k1}.
BIP 340, 2020. \url{https://github.com/bitcoin/bips/blob/master/bip-0340.mediawiki}.

\bibitem{taproot}
P. Wuille, J. Nick, and A. J. Towns. \emph{Taproot: SegWit Version 1
Spending Rules}. BIP 341, 2020.
\url{https://github.com/bitcoin/bips/blob/master/bip-0341.mediawiki}.

\bibitem{tls}
E. Rescorla. \emph{The Transport Layer Security (TLS) Protocol Version 1.3}.
RFC 8446, August 2018. \url{https://www.rfc-editor.org/rfc/rfc8446}.

\bibitem{tlsupdate}
E. Rescorla. \emph{The Transport Layer Security (TLS) Protocol Version 1.3}.
RFC 9846, 2026. Update superseding RFC 8446.
\url{https://www.rfc-editor.org/rfc/rfc9846}.

\end{thebibliography}
\end{document}
